\chapter{Euler telescoping and very-well-poised summations}
\label{chap:very-well-poised}

The summations in this chapter are finite.  Their natural first home is
therefore a rational-function field, where no convergence hypothesis is
needed and every nonzero auxiliary rational function may legitimately be
cancelled.  Only after the rational identity has been proved do we specialize
the parameters to complex numbers.  This order matters: normalizing a sum by
its conjectured product can introduce auxiliary denominators whose zeros are
not poles of the original identity.

The main result is Jackson's terminating balanced very-well-poised
${}_8\phi_7$ summation.  Its proof has two elementary layers.  A six-variable
rational transformation supplies the algebraic certificate, and Euler's
product-difference lemma turns that certificate into a finite telescoping
sum.  The terminal ${}_6\phi_5$ identity is then an honest finite limit.

\section{Euler's product-difference telescope}
\label{qg:sec:euler-telescope}

\begin{lemma}[Euler's telescoping lemma]
\label{qg:lem-euler-telescope}
Let $K$ be a field and let $N\in\NaturalNumbers$.  Choose
$u_0,\ldots,u_N\in K$ and $v_0,\ldots,v_N\in K$.  Put $w_k=u_k-v_k$ and assume
$w_0v_1\cdots v_N\ne0$.  Then
\begin{equation}
\label{qg:eq:euler-telescoping}
 \sum_{k=0}^{N}
 \frac{w_k}{w_0}
 \frac{u_0u_1\cdots u_{k-1}}{v_1v_2\cdots v_k}
 =\frac1{w_0}
 \left(
  \frac{u_0u_1\cdots u_N}{v_1v_2\cdots v_N}-v_0
 \right),
\end{equation}
where empty products are $1$.
\end{lemma}

\begin{proof}
For $k\ge1$ the $k$th summand is
\[
 \frac1{w_0}
 \left(
  \frac{u_0u_1\cdots u_k}{v_1v_2\cdots v_k}
  -\frac{u_0u_1\cdots u_{k-1}}{v_1v_2\cdots v_{k-1}}
 \right),
\]
because $w_k=u_k-v_k$.  The $k=0$ summand is
$(u_0-v_0)/w_0$.  Adding the $N+1$ differences cancels every intermediate
product and leaves exactly the right-hand side of
\eqref{qg:eq:euler-telescoping}.
\end{proof}

The application below has $u_N=v_0=0$, so the right-hand side vanishes.  The
advantage of \Cref{qg:lem-euler-telescope} is that the telescoping
antiderivative is encoded by the two product sequences; it need not be guessed
term by term.

\section{Jackson's rational certificate}
\label{qg:sec:jackson-certificate}

\begin{lemma}[Jackson's six-variable rational certificate]
\label{qg:lem-jackson-rational-certificate}
In the rational-function field
\[
 \RationalNumbers(A,B,C,D,E,F)
\]
one has
\begingroup
\scriptsize
\begin{align}
&1-
 \frac{(1-B)(1-C)(1-D)(1-E)(1-F)(1-A^3/(BCDEF))}
 {(1-A/B)(1-A/C)(1-A/D)(1-A/E)(1-A/F)(1-BCDEF/A^2)}
 \notag\\
&=\frac{(1-A)(1-D)(1-A^2/(BCDE))(1-A^2/(BCDF))
 (1-A^2/(BDEF))(1-A^2/(CDEF))}
 {(1-A/B)(1-A/C)(1-A/E)(1-A/F)
 (1-A^2/(BCDEF))(1-A^3/(BCD^2EF))}
 \notag\\
&\quad\times\left[
 1-
 \frac{(1-A/(BD))(1-A/(CD))(1-A/(DE))(1-A/(DF))
 (1-A^2/(BCDEF))(1-A^3/(BCDEF))}
 {(1-1/D)(1-A/D)(1-A^2/(BCDE))(1-A^2/(BCDF))
 (1-A^2/(BDEF))(1-A^2/(CDEF))}
 \right].
\label{qg:eq:jackson-rational}
\end{align}
\endgroup
Consequently the identity holds after every specialization to a field for
which all rational expressions displayed in
\eqref{qg:eq:jackson-rational} are defined.
\end{lemma}

\begin{proof}
For a finite list set
\[
 \Pi(x_1,\ldots,x_r):=\prod_{i=1}^{r}(1-x_i),
\]
and abbreviate the left-hand side of
\eqref{qg:eq:jackson-rational} by
\begin{align*}
J(A;B,C,D,E,F)
:={}&1-
\frac{\Pi(B,C,D,E,F,A^3/(BCDEF))}
{\Pi(A/B,A/C,A/D,A/E,A/F,BCDEF/A^2)}.
\end{align*}
Define also
\begin{align*}
\JacksonPrefactor(A;B,C,D,E,F)
:={}&
\frac{\Pi(A,A/(EF),A^2/(BCDE),A^2/(BCDF))}
{\Pi(A/E,A/F,A^2/(BCD),A^2/(BCDEF))}
\end{align*}
and the rational substitution
\begin{equation}
\label{qg:eq:jackson-seed-substitution}
\sigma(A;B,C,D,E,F)
:=\left(
 \frac{A^2}{BCD};
 E,F,\frac{A}{BC},\frac{A}{CD},\frac{A}{BD}
\right).
\end{equation}
We first prove the seed transformation
\begin{equation}
\label{qg:eq:jackson-seed}
 J(A;B,C,D,E,F)
 =\JacksonPrefactor(A;B,C,D,E,F)\,
  J\bigl(\sigma(A;B,C,D,E,F)\bigr).
\end{equation}

Let $D_0,N_0$ be respectively the denominator and numerator products in the
definition of $J$:
\begin{align*}
D_0&=\Pi(A/B,A/C,A/D,A/E,A/F,BCDEF/A^2),\\
N_0&=\Pi(B,C,D,E,F,A^3/(BCDEF)).
\end{align*}
After the substitution \eqref{qg:eq:jackson-seed-substitution}, the analogous
products, before common factors are cancelled, may be written
\begin{align*}
D_1&=\Pi(A/B,A/C,A/D,A^2/(BCDE),A^2/(BCDF),EF/A),\\
N_1&=\Pi(A/(BC),A/(BD),A/(CD),E,F,A^3/(BCDEF)).
\end{align*}
Put $\Delta_i=D_i-N_i$.  We claim
\begin{equation}
\label{qg:eq:jackson-delta-transform}
 \Delta_0=\lambda\Delta_1,
 \qquad
 \lambda:=\frac{B^2C^2D^2(A-1)}{A(A^2-BCD)}.
\end{equation}

This claim has a short coefficient proof.  Put $L=A^3/(BCDE)$ and regard
$\Delta_0,\Delta_1$ as Laurent polynomials in $F$.  Directly pairing the
constant and reciprocal factors shows that each has the form
\[
 \Delta_i=r_iF+s_i+\frac{Lr_i}{F}.
\]
Introduce the two four-term polynomials
\begin{align*}
U&:=A^2-A(B+C+D)+A+BCD,\\
V&:=A^2+ABCD-A(BC+BD+CD)+BCD.
\end{align*}
Expanding only in the affine variable $E$ gives
\begin{align*}
r_0&=\frac{A-1}{A^2}\,(-AU+EV),\\
r_1&=\frac{A^2-BCD}{AB^2C^2D^2}\,(-AU+EV).
\end{align*}
Hence $r_0=\lambda r_1$.

It remains to compare the constant coefficients.  At $F=1$ both $N_0$ and
$N_1$ vanish, and cancelling their common factors gives
\begin{align*}
\frac{\Delta_0(1)}{\Delta_1(1)}
&=\frac{(1-A/E)(1-A)(1-BCDE/A^2)}
{(1-A^2/(BCDE))(1-A^2/(BCD))(1-E/A)}\\
&=\frac{B^2C^2D^2(A-1)}{A(A^2-BCD)}
=\lambda.
\end{align*}
The coefficients of $F$ and $F^{-1}$ already have ratio $\lambda$, so the
equality at $F=1$ forces $s_0=\lambda s_1$.  This proves
\eqref{qg:eq:jackson-delta-transform}.  A second direct cancellation gives
\[
 \JacksonPrefactor(A;B,C,D,E,F)\frac{D_0}{D_1}=\lambda.
\]
Since $J=\Delta_0/D_0$ and
$J\circ\sigma=\Delta_1/D_1$, the seed transformation
\eqref{qg:eq:jackson-seed} follows.

Apply the seed transformation twice.  Its second iterate is
\begin{equation}
\label{qg:eq:jackson-second-iterate}
\sigma^2(A;B,C,D,E,F)
=\left(
 \frac{A^3}{BCD^2EF};
 \frac{A}{CD},\frac{A}{BD},\frac{A^2}{BCDEF},
 \frac{A}{DF},\frac{A}{DE}
\right).
\end{equation}
Substitution and cancellation in the two prefactors gives
\begin{align*}
&\JacksonPrefactor(A;B,C,D,E,F)\,
 \JacksonPrefactor\bigl(\sigma(A;B,C,D,E,F)\bigr)\\
&\quad=
\frac{\Pi(A,D,A^2/(BCDE),A^2/(BCDF),
 A^2/(BDEF),A^2/(CDEF))}
{\Pi(A/B,A/C,A/E,A/F,A^2/(BCDEF),A^3/(BCD^2EF))}.
\end{align*}
Likewise \eqref{qg:eq:jackson-second-iterate} gives
\begin{align*}
J(\sigma^2)
={}&1-
 \frac{(1-A/(BD))(1-A/(CD))(1-A/(DE))(1-A/(DF))}
 {(1-1/D)(1-A/D)(1-A^2/(BCDE))(1-A^2/(BCDF))}\\
&\qquad\times
 \frac{(1-A^2/(BCDEF))(1-A^3/(BCDEF))}
 {(1-A^2/(BDEF))(1-A^2/(CDEF))}.
\end{align*}
Thus
$J=\JacksonPrefactor\,(\JacksonPrefactor\circ\sigma)\,(J\circ\sigma^2)$
is precisely
\eqref{qg:eq:jackson-rational}.  Every step took place in the rational-function
field, so the final specialization assertion follows by applying a field
homomorphism wherever the displayed denominators remain nonzero.
\end{proof}

\section{Jackson's terminating \texorpdfstring{${}_8\phi_7$}{8phi7} sum}
\label{qg:sec:jackson-8phi7}

\begin{theorem}[Jackson's terminating very-well-poised ${}_8\phi_7$ sum]
\label{qg:thm-jackson-8phi7}
For $N\in\NaturalNumbers$, the following is an identity in
$\RationalNumbers(q,a,b,c,d)$:
\begin{align}
&\sum_{k=0}^{N}
 \frac{1-aq^{2k}}{1-a}
 \frac{(a,b,c,d,a^2q^{N+1}/(bcd),q^{-N};q)_k}
 {(q,aq/b,aq/c,aq/d,bcdq^{-N}/a,aq^{N+1};q)_k}q^k
 \notag\\
&\qquad=
 \frac{(aq,aq/(bc),aq/(bd),aq/(cd);q)_N}
 {(aq/b,aq/c,aq/d,aq/(bcd);q)_N}.
\label{qg:eq:jackson-8phi7}
\end{align}
Consequently \eqref{qg:eq:jackson-8phi7} holds for
$q,a,b,c,d\in\ComplexNumbers^\times$ whenever
\[
 (1-a)
 (q,aq/b,aq/c,aq/d,bcdq^{-N}/a,aq^{N+1},aq/(bcd);q)_N
 \ne0.
\]
Equivalently, it holds at every nonzero complex specialization at which the
terms in the displayed identity are defined.
\end{theorem}

\begin{proof}
Work first in the field $K=\RationalNumbers(q,a,b,c,d)$.  Let
$\JacksonNormalizedSummand{n}{k}$ be the
$k$th summand divided by the product on the right-hand side of
\eqref{qg:eq:jackson-8phi7}; explicitly,
\begin{align*}
\JacksonNormalizedSummand{n}{k}&=
 \frac{(aq/b,aq/c,aq/d,aq/(bcd);q)_n}
 {(aq,aq/(bc),aq/(bd),aq/(cd);q)_n}
 \frac{1-aq^{2k}}{1-a}q^k\\
&\quad\times
 \frac{(a,b,c,d,a^2q^{n+1}/(bcd),q^{-n};q)_k}
 {(q,aq/b,aq/c,aq/d,bcdq^{-n}/a,aq^{n+1};q)_k}.
\end{align*}
Set $\JacksonNormalizedSummand{n}{k}=0$ outside $0\le k\le n$.  Define
\begin{align*}
u_k={}&(1-aq^k)(1-bq^k)(1-cq^k)(1-dq^k)
 \left(1-\frac{a^2q^{n+k+1}}{bcd}\right)
 (1-q^{-n-1+k}),\\
v_k={}&\left(1-\frac{aq^k}{b}\right)
 \left(1-\frac{aq^k}{c}\right)
 \left(1-\frac{aq^k}{d}\right)
 \left(1-\frac{bcdq^{-n+k-1}}a\right)\\
&\times(1-aq^{n+k+1})(1-q^k),
\qquad
w_k:=u_k-v_k.
\end{align*}
Then $u_{n+1}=0$, $v_0=0$, and
\begin{equation}
\label{qg:eq:jackson-w0}
w_0=(1-a)(1-b)(1-c)(1-d)
 \left(1-\frac{a^2q^{n+1}}{bcd}\right)(1-q^{-n-1}).
\end{equation}

Apply \Cref{qg:lem-jackson-rational-certificate} with
\begin{equation}
\label{qg:eq:jackson-certificate-substitution}
A=aq^{2k},\quad B=\frac{aq^k}{b},\quad
C=\frac{aq^k}{c},\quad D=aq^{n+k+1},\quad
E=\frac{aq^k}{d},\quad F=q^k.
\end{equation}
The fraction subtracted from $1$ on the certificate's left-hand side becomes
$v_k/u_k$.  Cancelling the monomials on its right-hand side yields
\begin{equation}
\label{qg:eq:jackson-w}
w_k=bcd\,
 \frac{1-aq^{2k}}{1-a^2q^{2n+2}/(bcd)}
 \frac{q^k}{aq^{n+1}}\,\mathcal B_{n,k},
\end{equation}
where
\begin{align}
\mathcal B_{n,k}={}&
 (1-q^{-n-1})(1-aq^{n+1}/b)(1-aq^{n+1}/c)(1-aq^{n+1}/d)
 \notag\\
&\times(1-bcdq^{-n+k-1}/a)(1-a^2q^{n+k+1}/(bcd))
 \frac{aq^{n+1}}{bcd}
 \notag\\
&+(1-q^{-n+k-1})(1-aq^{n+k+1})
 (1-aq^{n+1}/(bc))(1-aq^{n+1}/(bd))
 \notag\\
&\times(1-aq^{n+1}/(cd))(1-a^2q^{n+1}/(bcd)).
\label{qg:eq:jackson-B}
\end{align}
The leading factor $bcd$ in \eqref{qg:eq:jackson-w} is essential.

We next verify the finite-difference certificate without an unprinted
multivariate expansion.  Put
\begin{align*}
\mathcal C_{n,k}:={}&-
 \frac{1-aq^{2k}}{1-a}q^k
 \frac{(a,b,c,d;q)_k}{(q,aq/b,aq/c,aq/d;q)_k}
 \frac{(aq/b,aq/c,aq/d,aq/(bcd);q)_n}
 {(aq,aq/(bc),aq/(bd),aq/(cd);q)_{n+1}}\\
&\times
 \frac{(q^{-n-1},a^2q^{n+1}/(bcd);q)_k}
 {(1-q^{-n-1})(1-a^2q^{n+1}/(bcd))
  (bcdq^{-n}/a,aq^{n+2};q)_k}.
\end{align*}
We claim that
\begin{equation}
\label{qg:eq:jackson-direct-difference}
 \JacksonNormalizedSummand{n+1}{k}-\JacksonNormalizedSummand{n}{k}
 =\mathcal C_{n,k}\mathcal B_{n,k}
 \qquad(0\le k\le n+1).
\end{equation}
For $0\le k\le n$, repeated use of
\[
 \frac{(xq;q)_k}{(x;q)_k}=\frac{1-xq^k}{1-x}
\]
gives
\begin{align*}
\frac{\JacksonNormalizedSummand{n+1}{k}}
     {\JacksonNormalizedSummand{n}{k}}-1
={}&-\frac{\mathcal B_{n,k}}
{(1-q^{-n+k-1})(1-aq^{n+1}/(bc))(1-aq^{n+1}/(bd))}\\
&\quad\times
\frac1{(1-aq^{n+1}/(cd))(1-a^2q^{n+1}/(bcd))
(1-aq^{n+k+1})}.
\end{align*}
Multiplying by the displayed formula for $\JacksonNormalizedSummand{n}{k}$
and cancelling one initial
and one terminal factor in each shifted factorial gives
\eqref{qg:eq:jackson-direct-difference}.  For $k=n+1$ one has
$(q^{-n};q)_{n+1}=0$, so $\JacksonNormalizedSummand{n}{n+1}=0$; the second summand of
\(\mathcal B_{n,n+1}\) also vanishes, and direct cancellation in the first
summand gives the same formula.

On the other hand,
\begin{equation}
\label{qg:eq:jackson-product-quotient}
 \frac{u_0u_1\cdots u_{k-1}}{v_1v_2\cdots v_k}
 =\frac{(a,b,c,d,a^2q^{n+1}/(bcd),q^{-n-1};q)_k}
 {(aq/b,aq/c,aq/d,bcdq^{-n}/a,aq^{n+2},q;q)_k}.
\end{equation}
Define
\begin{align*}
D_n={}&
 \frac{(aq/b,aq/c,aq/d,aq/(bcd);q)_n}
 {(aq,aq/(bc),aq/(bd),aq/(cd);q)_{n+1}}\\
&\times\left[-\frac{aq^{n+1}}{bcd}(1-b)(1-c)(1-d)
 \left(1-\frac{a^2q^{2n+2}}{bcd}\right)\right].
\end{align*}
The factor $bcd$ in the denominator of $D_n$ is essential.  Substituting
\eqref{qg:eq:jackson-w0}, \eqref{qg:eq:jackson-w}, and
\eqref{qg:eq:jackson-product-quotient}, and cancelling the common factors,
gives
\[
 D_n\frac{w_k}{w_0}
 \frac{u_0u_1\cdots u_{k-1}}{v_1v_2\cdots v_k}
 =\mathcal C_{n,k}\mathcal B_{n,k}.
\]
Together with \eqref{qg:eq:jackson-direct-difference}, this proves the exact
finite-difference certificate
\begin{equation}
\label{qg:eq:jackson-finite-difference}
 \JacksonNormalizedSummand{n+1}{k}-\JacksonNormalizedSummand{n}{k}
 =D_n\frac{w_k}{w_0}
 \frac{u_0u_1\cdots u_{k-1}}{v_1v_2\cdots v_k}.
\end{equation}

Use \Cref{qg:lem-euler-telescope} with upper index $n+1$.  Since
$u_{n+1}=v_0=0$, its right-hand side is zero.  Summing
\eqref{qg:eq:jackson-finite-difference} over $0\le k\le n+1$ therefore gives
\[
 \sum_k\JacksonNormalizedSummand{n+1}{k}
 =\sum_k\JacksonNormalizedSummand{n}{k}.
\]
At $n=0$ only $k=0$ occurs and
$\JacksonNormalizedSummand{0}{0}=1$.  Hence the normalized sum is $1$
for every $n$, which is exactly \eqref{qg:eq:jackson-8phi7} in $K$.

Finally, clear only the denominators occurring in the original identity.
Equality in $K$ then specializes at every complex tuple satisfying the
hypotheses of the theorem, even when an auxiliary normalization factor used
inside the proof vanishes.
\end{proof}

\begin{remark}[The removable value $a=1$]
\label{qg:rem:jackson-a-one}
The statement above excludes $a=1$ because the very-well-poised factor is
displayed as a quotient.  The singularity is removable.  For $k\ge1$,
\[
 \frac{1-aq^{2k}}{1-a}(a;q)_k
 =(1-aq^{2k})(aq;q)_{k-1},
\]
whose value at $a=1$ is $(1-q^{2k})(q;q)_{k-1}$; for $k=0$ the combined
factor is $1$.  Thus \eqref{qg:eq:jackson-8phi7} extends to $a=1$ whenever
the remaining displayed denominators are nonzero.
\end{remark}

\section{The terminating \texorpdfstring{${}_6\phi_5$}{6phi5} limit}
\label{qg:sec:terminating-6phi5}

\begin{corollary}[Terminating very-well-poised ${}_6\phi_5$]
\label{qg:cor-jackson-6phi5}
For $N\in\NaturalNumbers$, the identity
\begin{align}
&\sum_{k=0}^{N}
 \frac{1-aq^{2k}}{1-a}
 \frac{(a,b,c,q^{-N};q)_k}
 {(q,aq/b,aq/c,aq^{N+1};q)_k}
 \left(\frac{aq^{N+1}}{bc}\right)^k
 \notag\\
&\qquad=
 \frac{(aq,aq/(bc);q)_N}{(aq/b,aq/c;q)_N}
\label{qg:eq:terminating-6phi5}
\end{align}
holds in $\RationalNumbers(q,a,b,c)$.  Consequently it holds for
$q,a,b,c\in\ComplexNumbers^\times$ whenever
\[
 (1-a)(q,aq/b,aq/c,aq^{N+1};q)_N\ne0.
\]
It also has the removable $a=1$ extension described in
\Cref{qg:rem:jackson-a-one}.
\end{corollary}

\begin{proof}
In \eqref{qg:eq:jackson-8phi7}, regard $x=d^{-1}$ as an indeterminate over
$\RationalNumbers(q,a,b,c)$ and take the value at $x=0$ after cancelling the
common leading powers of $d$.  For each fixed $k$,
\begin{align*}
(d;q)_k
 &=(-d)^kq^{\binom{k}{2}}\bigl(1+O(d^{-1})\bigr),\\
(bcdq^{-N}/a;q)_k
 &=\left(-\frac{bcdq^{-N}}a\right)^k
 q^{\binom{k}{2}}\bigl(1+O(d^{-1})\bigr).
\end{align*}
Therefore
\[
 q^k\frac{(d;q)_k}{(bcdq^{-N}/a;q)_k}
 \longrightarrow
 \left(\frac{aq^{N+1}}{bc}\right)^k.
\]
The factors with arguments $a^2q^{N+1}/(bcd)$ and $aq/d$ tend to $1$.
On the product side the four factors involving $d$ likewise tend to $1$,
leaving the quotient in \eqref{qg:eq:terminating-6phi5}.  The sum has only
$N+1$ terms, so taking this coefficientwise rational limit is legitimate and
proves the identity in $\RationalNumbers(q,a,b,c)$.  The specialization and
removable-value statements follow exactly as in
\Cref{qg:thm-jackson-8phi7,qg:rem:jackson-a-one}.
\end{proof}
